\documentclass[a4paper,11pt,oneside]{article}

\makeatletter
\def\input@path{{../}{}}
\makeatother
\usepackage{style/project}

% Exercise Sheet 02

\begin{document}

\begin{center}
    {\LARGE\bfseries Groups and Lie Algebras: Exercise Sheet 2\par}
    \vspace{0.5em}
    {\color{gray!70}\rule{0.72\textwidth}{0.5pt}\par}
\end{center}

\setcounter{section}{2}

    \ecs{Simple groups}{
        A group \(G\) is called \textbf{simple} if it has exactly two normal
        subgroups: the trivial subgroup \(\{1\}\) and \(G\) itself.
    }{
        \ecpart{Let \(G\) be a simple group and \(H\) a group. Show that every
        group homomorphism from \(G\) to \(H\) is either an injection or the
        constant map at the identity of \(H\). Show that any group \(G'\)
        satisfying this property must be simple.}{
            Let \(f:G \to H \) be a group homomorphism. Then by the isomophism theorem the kernel if \(f\) is a normal subgroup of \(G\), i.e. \(\ker f \in \{\{1\}, G\}\). If \(\ker f = G\) then \(f\) is the constant map at the identity of \(H\). 
            To show that if \(\ker f = \{1\}\) impliels \(f\) is an injection consider \(g_1, g_2 \in G\) such that \(f(g_1)=f(g_2)\). Then \(f(g_2^{-1}g_1)=f(1) = 1\) and thus \(g_2^{-1}g_1 \in \ker f = \{1\}\) which implies \(g_1=g_2\).
            
            Conversely, let \(G^\prime\) be a group such that every group homomorphism \(h:G^\prime \to H\) for any group \(H\) is either an injection or the constant map at the identity of \(H\). Let \(N \leq G^\prime\) be a normal subgroup. Then there exists a group homomorphism \(f_1\) such that \(\ker f = N\). If \(f_1\) is an injection then \(\ker f = \{1\}\), and if \(f_1\) is the constant map at the identity of \(H\) then \(\ker f = G^\prime\). In either case \(N \in \{\{1\}, G^\prime\}\) and thus \(G^\prime\) is simple.
            }

        \ecpart{For which \(n \in \mathbb{N}\) is
        \(\mathbb{Z}/n\mathbb{Z}\) simple?}{
            Let \(N \leq \mathbb{Z}/n\mathbb{Z}\) be a normal subgroup. Since \(\mathbb{Z}/n\mathbb{Z}\) is abelian, every subgroup is normal. Since every subgroup of \(\mathbb{Z}/n\mathbb{Z}\) is of the form \(k\mathbb{Z}/n\mathbb{Z}\) for some \(k \in \mathbb{N}\) dividing \(n\). Therefore, \(\mathbb{Z}/n\mathbb{Z}\) is simple if and only if \(n\) has no divisors other than \(1\) and \(n\), i.e. if and only if \(n\) is prime.
        }

        \ecpart{Let \(G\) be a group and \(N \subset G\) a maximal normal
        subgroup, i.e. \(N \neq G\) and
        \[
            N \subsetneq N' \subset G
        \]
        with \(N'\) normal implies \(N' = G\). Show that \(G/N\) is simple.}{
            Consider the canonical projection \(\pi: G \to G/N, \quad g \mapsto [g]\) and \(H \trianglelefteq G/N\). Then the preimage \(\pi^{-1}(H)\) is a normal subgroup of \(G\) since for \(g \in G\) and \(h \in \pi^{-1}(H)\) we have \(\pi(ghg^{-1}) = \pi(g)\pi(h)\pi(g^{-1}) \in H\) since H is normal. But \(N \subset \pi^{-1}(H)\) since \(\pi^{-1}(1) = N\). Thus either \(\pi^{-1}(H) = N\) or \(\pi^{-1}(H) = G\) by the maximality of \(N\). In the first case \(H = 1\) and in the second case \(H = G/N\). Therefore, \(G/N\) is simple.
       }

        \ecpart{Show that a group with only two conjugacy classes is simple.}{
            First we show that every normal subgroup \( N \trianglelefteq G\) is a union of conjugacy classes. Let \(h \in N\), since \(N\) is normal we have \(ghg^{-1} \in N \, \forall g\in G\), i.e. the conjugacy class of \(h\) is contained in \(N\). Therefore, \(N\) is the union of the conjugacy classes of its elements.
            
            For every group \([1] = \{1\} \in \operatorname{conj}(G)\). If \(G\) has only two conjugacy classes, then the second conjugacy class is \(G/\{1\}\). Every normal subgroup \(N \trianglelefteq G\) is a union of conjugacy classes, so either \(N \in \{\{1\}, G/\{1\}, G\}\). But \(G/\{1\}\) is not a subgroup of \(G\) since it does not contain the identity, so \(N \in \{\{1\}, G\}\) and thus \(G\) is simple.
        }

        \ecpart{Iterating the construction from c), we can assign to each finite
        group \(G\) a sequence of simple groups. Show explicitly for
        \(G = \mathbb{Z}/12\mathbb{Z}\) that the simple groups appearing in this
        decomposition are independent of the choice of maximal normal subgroup
        (this is true for every finite group).}{
            To iterate the construction we start with \(G \trianglerighteq G_1\) where \(G_1\) is a maximal normal subgroup of \(G\). Then \(G/G_1\) is simple and therefore of prime order. This is iterated by choosing a maximal normal subgroup \(G_2\) of \(G_1\) and so on until we reach the trivial group. The simple groups appearing in this decomposition are the quotients \(G/G_1, G_1/G_2, \ldots\). We show that for \(G = \mathbb{Z}/12\mathbb{Z}\) the simple groups appearing in this decomposition are independent of the choice of maximal normal subgroup.
            Since \(G\) is abelian, every subgroup is normal. For \(\mathbb{Z}/k\mathbb{Z}\) to be a subgroup of \(\mathbb{Z}/12\mathbb{Z}\) we need \(k\) to divide \(12\). The divisors of \(12\) are \(1, 2, 3, 4, 6, 12\). For a subgroup to be maximal \(k\) needs to be a proper divisor of \(12\) such that there exists no proper divisor of \(12\) which is a multiple of \(k\). The maximal normal subgroups of \(\mathbb{Z}/12\mathbb{Z}\) are therefore \(\mathbb{Z}/4\mathbb{Z}\) and \(\mathbb{Z}/6\mathbb{Z}\).

            Starting with \(G_1 = \mathbb{Z}/4\mathbb{Z}\) we have \(G/G_1 \cong \mathbb{Z}/3\mathbb{Z}\). The maximal proper divisor of \(4\) is \(2\) and thus \(G_2 = \mathbb{Z}/2\mathbb{Z}\) and therefore the second simple group is \(G_1/G_2 \cong \mathbb{Z}/2\mathbb{Z}\). Finally, the maximal proper divisor of \(2\) is \(1\) and thus \(G_3 = \{0\}\) and the third simple group is \(G_2/G_3 \cong \mathbb{Z}/2\mathbb{Z}\). Therefore the decomposition yields the simple groups \(\mathbb{Z}/3\mathbb{Z}, \mathbb{Z}/2\mathbb{Z}, \mathbb{Z}/2\mathbb{Z}\).

            On the other hand, starting with \(G_1 = \mathbb{Z}/6\mathbb{Z}\) we have \(G/G_1 \cong \mathbb{Z}/2\mathbb{Z}\). The maximal proper divisors of \(6\) are \(2,3\) and thus \(G_2 = \mathbb{Z}/3\mathbb{Z}\) or \(G_2 = \mathbb{Z}/2\mathbb{Z}\). Therefore the second simple group is \(G_1/G_2 \cong \mathbb{Z}/2\mathbb{Z}\) or \(G_1/G_2 \cong \mathbb{Z}/3\mathbb{Z}\). Finally, the maximal proper divisor of \(3\) is \(1\) as it is of \(2\) and thus \(G_3 = \{0\}\) for both cases and the third simple group is \(G_2/G_3 \cong \mathbb{Z}/3\mathbb{Z}\) or \(G_2/G_3 \cong \mathbb{Z}/2\mathbb{Z}\). Therefore the decomposition yields the simple groups \(\mathbb{Z}/2\mathbb{Z}, \mathbb{Z}/2\mathbb{Z}, \mathbb{Z}/3\mathbb{Z}\).

            Thus all possible decompositions of \(\mathbb{Z}/12\mathbb{Z}\) yield the same simple groups \(\mathbb{Z}/3\mathbb{Z}, \mathbb{Z}/2\mathbb{Z}, \mathbb{Z}/2\mathbb{Z}\).
        }

        \ecpart{The decomposition in e) is called the Jordan-H\"older
        decomposition of \(G\). Show that there are non-isomorphic groups with
        the same Jordan-H\"older decomposition.}{

        }
    }

    \ec{Spontaneous symmetry breaking}{
        It can happen that a quantum mechanical system is in a ground state
        which does not preserve the symmetries of the system. In this case we
        say the symmetry is \textit{spontaneously broken}.

        Consider a particle on a circle with Hilbert space
        \[
            \mathcal{H} = L^2(\mathbb{R}/\mathbb{Z})
        \]
        and Hamiltonian
        \[
            \mathbb{H}
            =
            \frac{1}{2}\left(-i\partial_x - \frac{1}{2}\right)^2.
        \]
        Compute the eigenfunctions of \(\mathbb{H}\) and show that
        \[
            f(x) \mapsto f(-x)
        \]
        is a symmetry. What is the action of the symmetry when applied to the
        eigenfunctions?
    }{
        To find the eigenfunctions of \(\mathbb{H}\) we find the characteristic polynomial of \(\mathbb{H}\). For \(f \in \mathcal{H}\) we have
        \[
            \mathbb{H}f = \frac{1}{2}\left(-i\partial_x - \frac{1}{2}\right)^2f = \frac{1}{2}\left(-\partial_x^2 + i\partial_x + \frac{1}{4}\right)f = E f
        \]
        therefore the characteristic polynomial of \(\mathbb{H}\) with the ansatz \(f(x) = e^{\lambda x}\) is given by
        \[
            -\lambda^2 + i\lambda + \frac{1}{4} - 2E = 0
        \]
        The solutions of this polynomial are given by the pq formula
        \[
            \lambda = \frac i2 \pm \sqrt{-2E} = i\left(\frac 12 \pm \sqrt{2E}\right)
        \]
        and thus the eigenfunctions of \(\mathbb{H}\) are given by
        \[
            f(x) = Ae^{i\left(\frac 12 + \sqrt{2E}\right)x} + Be^{i\left(\frac 12 - \sqrt{2E}\right)x}
        \]
        for some \(A, B \in \mathbb{C}\). Since \(f\) is takes values on the circle, we have \(f(x) = f(x+2\pi)\) and thus
        \begin{align*}
            Ae^{i\left(\frac 12 + \sqrt{2E}\right)x} + Be^{i\left(\frac 12 - \sqrt{2E}\right)x} &= Ae^{i\left(\frac 12 + \sqrt{2E}\right)(x+2\pi)} + Be^{i\left(\frac 12 - \sqrt{2E}\right)(x+2\pi)} \\
            &= Ae^{i\left(\frac 12 + \sqrt{2E}\right)x}e^{i\left(\frac 12 + \sqrt{2E}\right)2\pi} + Be^{i\left(\frac 12 - \sqrt{2E}\right)x}e^{i\left(\frac 12 - \sqrt{2E}\right)2\pi} \\
        \end{align*}
        and thus
        \[
            e^{i\left(\frac 12 + \sqrt{2E}\right)2\pi} = e^{i\left(\frac 12 - \sqrt{2E}\right)2\pi} = 1
        \]
        which implies \(\frac12 \pm \sqrt{2E} \in \mathbb{Z}\) and thus \(E = \frac{1}{2}\left(n \mp \frac{1}{2}\right)^2\) for some \(n \in \mathbb{Z}\). Therefore, the eigenfunctions of \(\mathbb{H}\) are given by
        \[
            f(x) = Ae^{inx} + Be^{-inx}
        \]
        for some \(A, B \in \mathbb{C}\) and \(n \in \mathbb{Z}\). The eigenfunctions of \(\mathbb{H}\) transform under parity as
        \[
            f(-x) = Ae^{-inx} + Be^{inx}
        \]
        which is still an eigenfunction of \(\mathbb{H}\) with the same eigenvalue \(E\). Therefore, the parity transformation is a symmetry of the system. 
    }

    \ecs{Symmetries of the Dodecahedron}{
        In this exercise we study the symmetry group of the dodecahedron.
    }{
        \ecpart{Compute the size of the group of symmetries of the dodecahedron
        using the orbit-stabilizer theorem.}{
            Concider the faces of the dodecahedron. The orbit of a face are all the faces of the dodecahedron. Therefore \(\mid \mathcal{O}(F)\mid = 12\).
            The stabilizer of a face is the symmetries of a pentagon, which consists of \(5\) rotations and \(5\) reflections, and thus \(\mid \operatorname{Stab}(F)\mid = 10\).
            Therefore, by the orbit-stabilizer theorem, the size of the group of symmetries of the dodecahedron is given by \(\mid \operatorname{Sym}(\text{dodecahedron})\mid = \mid \mathcal{O}(F)\mid \cdot \mid \operatorname{Stab}(F)\mid = 12 \cdot 10 = 120\).
        }

        \ecpart{What is the size of the group of rotational symmetries?}{
            The group of rotations is the subgroup which does not contain reflections. Since the group of symmetries of the dodecahedron contains \(5\) reflections for each face, the group of rotations contains \(5\) rotations for each face. Therefore, the size of the group of rotational symmetries is given by \(\mid \operatorname{Rot}(\text{dodecahedron})\mid = 12 \cdot 5 = 60\).
        }

        \ecpart{There are \(5\) cubes inside the dodecahedron which can be
        rotated into each other and contain all its vertices. Show that there is
        an injection from the group of rotational symmetries to \(S_5\).}{
            To inscribe a cube into a dodecahedron, notice that the number of edges on a cube and the number of faces on a dodecahedron is 12. A cube is uniquely defined by a diagonal on a face of the dodecahedron. Since each face has 5 vertices there are 5 distinct cubes. Laber this cubes by \(\mathcal{C} = \{\mathcal{C}_1, \cdots , \mathcal{C}_5\}\). The group of rotation preserves the order of vertices and thus the cubes. Therefore, there is a group action of the rotational group on the set of cubes:
            \[
                \operatorname{Rot}_{12} \times \mathcal{C} \to \mathcal{C}, \quad (r, \mathcal{C}_i) \mapsto r\triangleright \mathcal{C}_i
            \]
            This induces a bijection \(\sigma(g) : \mathcal{C} \to \mathcal{C}\) which is isomorphic to a permutation of the set \(\{1, \cdots , 5\}\). Define the map 
            \[\phi : \operatorname{Rot}_{12} \to S_5, \quad g \mapsto \sigma(g)\]
            \(\phi\) is a homomorphism since for \(g_1, g_2 \in \operatorname{Rot}_{12}\) and \(\mathcal{C}_i \in \mathcal{C}\) we have
            \begin{align*}
                \phi(g_1g_2)(\mathcal{C}_i) &= \sigma(g_1g_2)(\mathcal{C}_i) \\
                &= (g_1g_2)\triangleright \mathcal{C}_i \\
                &= \sigma(g_1)(\sigma(g_2)(\mathcal{C}_i)) \\
                &= (\sigma(g_1)\circ \sigma(g_2))(\mathcal{C}_i) \\
                &= (\phi(g_1)\circ \phi(g_2))(\mathcal{C}_i)
            \end{align*}
            \(\phi\) is injective since for \(g_1, g_2 \in \operatorname{Rot}_{12}\) such that \(\phi(g_1) = \phi(g_2)\) we have:
            \begin{align*}
                \phi(g_1)(\mathcal{C}_i) &= \phi(g_2)(\mathcal{C}_i) \\
                \sigma(g_1)(\mathcal{C}_i) &= \sigma(g_2)(\mathcal{C}_i) \\
                g_1\triangleright \mathcal{C}_i &= g_2\triangleright \mathcal{C}_i \\
                g_2^{-1}g_1\triangleright \mathcal{C}_i &= \mathcal{C}_i
            \end{align*}
            and therefore \(g_2^{-1}g_1\) is the identity. Thus \(g_1 = g_2\) and \(\phi\) is injective.
        }

        \ecpart{Show that the image of the injection from c) is contained in the
        alternating group \(A_5\).}{
            
        }

        \ecpart{Show that the group of symmetries of the dodecahedron is
        \(A_5 \times \mathbb{Z}_2\).}{

        }
    }

    \ec{Cayley's theorem}{
        Show that every finite group is isomorphic to a subgroup of a symmetric
        group.

        Hint: Use the action of a group on itself by left multiplication.
    }{
        Let \(G\) be a finite group. The action of \(G\) on itself by left multiplication is given by \(L_g[h] \equiv g \triangleright  h = gh \). Define the map \(\phi: G \to \operatorname{Sym}_G, \quad g \mapsto L_g\).

        First we show that this is well defined, i.e. that \(L_g\) is a bijective homomorphism for every \(g \in G\). 
        \begin{itemize}
            \item \(L_g\) is a homomorphism: For \(h_1, h_2 \in G\) we have \(L_g[h_1h_2] = g(h_1h_2) = (gh_1)(gh_2) = L_g[h_1]L_g[h_2]\).
            \item \(L_g\) is injective: For \(h_1, h_2 \in G\) such that \(L_g[h_1] = L_g[h_2]\) we have \(gh_1 = gh_2\) and thus \(h_1 = h_2\).
            \item \(L_g\) is surjective: For \(h \in G\) we have \(L_g[g^{-1}h] = g(g^{-1}h) = h\).
        \end{itemize}

        Next we show that \(\phi\) is a homomorphism. For \(g_1, g_2, h \in G\) we have \[\phi(g_1g_2)[h] = L_{g_1g_2}[h] = (g_1g_2)h = g_1(g_2h) = L_{g_1}[L_{g_2}[h]] = (\phi(g_1)\circ \phi(g_2))[h]\]

        Finally to show that \(\phi\) is an isomorphism we show that it is injective. For \(g_1, g_2 \in G\) such that \(\phi(g_1) = \phi(g_2)\) we have \(L_{g_1} = L_{g_2}\) and thus \(g_1h = g_2h\) for every \(h \in G\). In particular, for \(h = 1\) we have \(g_1 = g_2\). Therefore, \(\phi\) is an injection and thus \(G\) is isomorphic to a subgroup of \(\operatorname{Sym}_G\).
    }

\end{document}
